\section{ARMA 与其他方法的连接：什么时候用什么}
\label{sec:06-Random-Processes/06-ARMA-and-ARIMA/04}

一个需求预测的项目落到你手上：五千个 SKU，每个有三年的周度销量，也就是每条序列一百五十个点左右。会上分成两派，一派要上序列到序列的深度模型，一派说 ARIMA 跑一遍就够。

两派各自都能举出证据，因为他们在回答不同的问题。一百五十个点用来估计一个 ARIMA(1,1,1) 绰绰有余，用来训练一个几十万参数的网络显然不够；可这里有五千条序列，如果它们共享同一套季节形状和促销响应，那就不是一百五十个点，是七十五万个点。选型的分歧往往不在算法优劣，而在于对``我到底有多少数据''这个问题的两种读法。

\subsection{先把坐标定下来}

值得先摆清楚两条轴：一条是可用信息量，另一条是数据偏离 ARMA 假设的程度。ARMA 的假设很具体——线性、平稳（或差分后平稳）、参数不随时间变、创新方差恒定、单序列自足。任何一条被违反的严重程度，都可以往上面那条轴上放。

\begin{center}
\begin{tikzpicture}[x=1cm,y=1cm,font=\small,>=stealth]
  \fill[blue!7]  (0,0) rectangle (3.4,2.1);
  \fill[green!8] (0,2.1) rectangle (3.4,4.7);
  \fill[red!7]   (3.4,2.1) rectangle (8.2,4.7);
  \fill[blue!4]  (3.4,0) rectangle (8.2,2.1);
  \draw[black!35] (0,0) rectangle (8.2,4.7);
  \draw[black!35] (3.4,0) -- (3.4,4.7);
  \draw[black!35] (0,2.1) -- (8.2,2.1);
  \draw[->,black!60] (0,-0.15) -- (8.4,-0.15) node[right,font=\footnotesize,align=left] {可用信息量\\\footnotesize（序列长度 \(\times\) 序列条数）};
  \draw[->,black!60] (-0.15,0) -- (-0.15,4.9) node[above,font=\footnotesize] {偏离 ARMA 假设的程度};
  \node[font=\footnotesize,align=center,black!80] at (1.7,1.05) {ARMA / ARIMA\\\footnotesize 参数少、可解释\\\footnotesize 短期预测难被打败};
  \node[font=\footnotesize,align=center,black!80] at (1.7,3.4) {结构化扩展\\\footnotesize 状态空间、门限模型\\\footnotesize GARCH、季节分解};
  \node[font=\footnotesize,align=center,black!80] at (5.8,3.4) {跨序列的学习模型\\\footnotesize 梯度提升、深度序列模型\\\footnotesize 靠共享参数换容量};
  \node[font=\footnotesize,align=center,black!80] at (5.8,1.05) {ARMA 加外生变量\\\footnotesize 数据多但结构简单时\\\footnotesize 复杂模型换不来增益};
  \draw[black!45,dashed] (0.9,4.62) -- (8.05,1.05);
\end{tikzpicture}

\emph{两条轴，四个落点。虚线是预算线：模型容量越大，需要的信息量越多，越过这条线往左上走就是过拟合，往右下走则是把数据浪费掉。注意右下角——数据很多但结构简单时，复杂模型带不来增益，这一格常被忽略。}
\end{center}

图里最容易被误读的是横轴。它不是``行数''，是可用于估计同一批参数的信息量。五千条独立建模的 ARIMA，每条只用自己那一百五十个点，横轴位置在左；一个在五千条序列上共享权重的模型，横轴位置在右。同一份数据，两种建模方式落在图的两端。深度模型在实际业务预测里赢下来的场合，多半是靠这次横移，而不是靠非线性。

\subsection{状态空间不是 ARMA 的对手，是它的上位框架}

先处理一个常见的误解：ARMA 和状态空间不是两条并行路线。ARMA(\(p,q\)) 可以精确地写成一个线性 Gaussian 状态空间模型，它的最小均方误差预测与稳态 Kalman 滤波给出的预测完全一致（见\hyperref[sec:13-Machine-Learning/07-Graphical-and-Sequential-Models/03]{``Kalman filter 是动态系统上的 Gaussian 条件化''}）。选择哪一边，是选表述方式，不是选模型类。

那什么时候值得付出改写的成本？看四种症状。观测里有缺失或者采样间隔不齐时，ARMA 的似然要专门打补丁，状态空间只需在缺失那一步跳过更新。参数随时间缓慢演化时，状态空间可以把参数本身当成状态，让它随数据漂；ARMA 只能滚动窗口重估。有多个观测通道测同一个隐藏量时——三个传感器测同一个温度、多个指标反映同一个健康度——状态空间天然表达这种结构。需要把不确定性拆成``隐藏状态没看清''和``观测有噪声''两部分时，只有状态空间给得出这个分解。

反过来，序列平稳、参数稳定、观测规整时，多写一层状态方程只会增加实现负担与识别难度。ARMA 的诊断工具成熟得多，两条相关图加一个 Ljung--Box 就能把大部分问题暴露出来。

\subsection{机器学习买到什么，付出什么}

把滞后值当特征扔进梯度提升或神经网络，这条路的三项真实收益值得分开说。

最大的一项是跨序列共享。五千个 SKU 各自的历史都短，但促销的形状、周内的节律、生命周期的形态是共通的；一个在全部序列上训练的模型可以从别的 SKU 身上学到这些，再迁移到只有二十周历史的新品上。逐序列的 ARIMA 做不到这件事——它连``别的序列存在''都不知道。冷启动和长尾序列是这类模型压倒性占优的地方。

第二项是外生特征。价格、促销标记、天气、节假日、竞品动作，这些变量与目标的关系常常是非线性且有交互的（打折深度对销量的作用在节假日和平日完全不同）。ARMA 家族要靠人工设定形式才能吸收这些，学习模型可以自己找。

第三项才是非线性本身，而它通常被高估。很多序列在做完对数变换和季节调整之后，剩下的依赖结构接近线性；此时非线性容量带来的增益很小，噪声却被更充分地记住了。

代价同样有三项。样本外区间估计要靠分位数回归或保形预测另外做一遍，不像 ARIMA 那样从模型直接读出来。诊断手段稀薄，模型失效时不容易定位到具体哪条假设破了。最容易吃亏的是外推：树模型的预测是训练集里叶子均值的组合，训练时没见过的取值范围它一律给不出来，一条持续上行的序列直接喂给梯度提升，预测会在最后一个水平上压平。所以差分或去趋势这一步在树模型里不是可选项，它比在 ARIMA 里更关键——上一节讲的那套判断在这里照样成立，甚至更严格。

选型的实际经验大致是这样：单条长序列、结构简单、需要区间和解释，经典方法通常更稳；成千上万条短序列、有丰富协变量、评价指标只看点预测误差，学习模型更有机会。公开预测竞赛的结果也支持这个分界——单序列赛道上纯统计方法长期难以撼动，直到把指数平滑与循环网络混合起来才被超过；而在有大量商品、门店与促销信息的零售赛道上，梯度提升配上精心构造的特征几乎是标配。

\subsection{频域看到的是同一个模型的另一面}

ARMA 与谱分析也不是两套东西。平稳过程的谱密度与自协方差互为 Fourier 变换，ARMA 的谱密度可以直接写出来：

\[
f(\omega)=\frac{\sigma^2}{2\pi}
\left\lvert\frac{\theta(e^{-i\omega})}{\phi(e^{-i\omega})}\right\rvert^2 .
\]

分母的根靠近单位圆时，对应频率上出现尖峰；分子的根靠近单位圆时，对应频率被压成谷。上一节说差分是高通滤波器，用的就是这套语言。所以 AR 的伪周期、MA 的凹陷、差分的代价，在频域里都是同一张图上的几何特征。

值得换到频域的场合是周期本身就是研究对象的时候：有多个不同尺度的周期叠在一起、需要设计滤波器保留或去掉某个频段、或者要判断某个可疑的周期成分是真实的还是采样造成的混叠。反过来，如果最终交付的是下三期的点预测和区间，时域模型直接得多——谱估计给不出预测公式。

\subsection{三种情况下 ARMA 应该直接放弃}

有些结构不是调阶数能救的。

一是状态切换。系统在几种机制之间跳，每种机制下的动力学不同：拥塞与非拥塞、牛市与熊市、正常与故障。ARMA 会把两套动力学平均成一套，得到的参数在任何一个机制里都不对。门限自回归按可观测的阈值切换，Markov 切换模型让机制本身是一条不可观测的 Markov 链，后者与\hyperref[sec:06-Random-Processes/03-Markov-Chains/01]{``Markov 性质、转移矩阵与多步演化''}是同一套工具。

二是多变量强耦合。几条序列互为因果时，逐条建 ARMA 会丢掉全部交叉信息。向量自回归是直接的推广，但要留意代价：\(K\) 条序列 \(p\) 阶滞后有 \(K^2p\) 个参数，十条序列四阶就是四百个，样本很快不够用。这时要么加正则化，要么降维再建模。若各序列都有单位根而某个线性组合平稳，则属于协整的范畴，下一节会碰到。

三是方差本身有动态。均值结构可能已经很弱了——金融收益率的自相关常常接近零——真正有结构的是波动率。这时 ARMA 拟合出的残差是白的，残差平方却不白，说明该建模的对象换了一个。见\hyperref[sec:11-Mathematical-Statistics/10-GARCH-and-Volatility/01]{``条件异方差的统计特征''}。ARMA 管均值、GARCH 管方差，这不是谁替代谁，是分工。

\subsection{选型之前先立一条基线}

上面所有比较都建立在一个前提上：你有办法判断谁更好。这件事经常做得比选模型还草率。

样本内的信息准则不足以决定选型。AIC 挑的是预测导向的模型，BIC 挑的是更接近真模型的模型，两者在同一份数据上常给出不同答案；更麻烦的是，它们都建立在模型设定正确这个假设上，而选型这个动作恰恰是在设定还不确定的时候做的。可靠的做法是滚动样本外评估：按时间切分，只用切分点之前的数据训练，预测之后若干期，再往前挪一格重复。任何跨越时间边界的操作——用全样本做标准化、用未来数据填缺失、按随机而非时间划分验证集——都会让评估结果偏乐观，而且偏得很难察觉。

同样重要的是立基线。随机游走（明天等于今天）和季节朴素（今天等于上一个周期的今天）几乎不需要成本，却经常打败调了两周参数的模型。任何新方案先和它们比，比不过就说明信号不在你以为的地方。这条纪律看起来朴素，但它拦下的错误比任何模型选择准则都多。

到这里，讨论的一直是序列自己解释自己。可现实里你常常还握着别的变量——价格、投放、上游库存、气温——它们对目标有影响，而且影响不是当期就结清的。下一节把这些变量接进来，看回归和 ARMA 这两套语言怎样合成一个动态模型。
